\begin{papercard}{huckvale2015}
\cf{Stages} \stagetag{A} \stagetag{D} \stagetag{R} (also \stagetag{C}, \stagetag{W}). A systematic six-month assessment of all 79 apps in an accredited library (NHS Health Apps Library). Simulated-use testing, local storage inspection and traffic interception were compared with an ICO-guidance policy coding schema.

\cf{Aim} Assess whether accredited apps comply with UK data protection principles: collection, transmission and local storage; confidentiality; availability and content of policies; and agreement between policy and behaviour \pc{huckvale2015}{2}{The purpose of the current study was to assess the ex}.

\cf{Sampling frame} \emph{All} apps in the library in July 2013 for Android or iOS, a complete census of a curated list. Exclusions: not downloadable; costing $>$50 USD; free/demo/lite versions when a full version exists; not starting after two attempts on different devices. Cross-platform apps were tested separately but combined for analysis. No app or OS updates were applied during evaluation \pc{huckvale2015}{2}{All mobile apps available in the NHS Health Apps}. 88 apps, 9 excluded, 79 included \pc{huckvale2015}{5}{In July 2013, the NHS Health Apps Library contained 88}.

\cf{Simulated-use protocol} Six months (August 2013 to January 2014). First, sequential exploration of \emph{all} UI elements for coverage; then extended routine and intermittent use to catch infrequent behaviours such as delayed feedback requests or batched analytics \pc{huckvale2015}{3}{Testing incorporated}. Two researchers worked independently with simulated, clinically plausible data and app-specific accounts \pc{huckvale2015}{3}{Testing was performed by two study researchers}. Ethics device: test messages sent through support or feedback channels were flagged as tests, and this was done only at the end to avoid biasing data handling \pc{huckvale2015}{3}{To avoid misleading poten}. No human participants, so no consent was needed.

\cf{Data-type schema} Entered and transmitted data were coded with a schema (AF1) of strong identifiers, weaker identifiers, health-related information, other sensitive information, other user-generated data and protocol-standard data \pc{huckvale2015af}{1}{Data types identified in app data collections and transmissions}.

\cf{Storage and traffic} Device files and caches were copied and inspected for content and encryption \pc{huckvale2015}{3}{At the conclusion of the evaluation period, data stored on each test}. Traffic was captured by a man-in-the-middle set-up that also decrypted TLS, piloted beforehand on non-study apps \pc{huckvale2015}{3}{To capture data transmitted by included apps}. Two-step coding: automatic tagging by destination, security and known test values, then manual review by two independent researchers, reconciled by discussion \pc{huckvale2015}{3}{Personal information sent by apps was categorized in a}. Server-side checks: replaying requests without session state, and one confirmed SQL injection, without modifying or deleting data \pc{huckvale2015}{4}{Potential vulnerabilities in server-side controls were}.

\cf{Policy review} Sources searched: in-app documentation, store entries, related websites. Every document self-identifying as a privacy policy was included, plus privacy-relevant text in terms, disclaimers and store marketing text \pc{huckvale2015}{4}{We searched in-app documentation, app store entries and related websites}. The coding schema was based on UK ICO guidance under the Data Protection Act 1998 (AF4: uses of data, technical concerns, user rights, administrative details) \pc{huckvale2015}{4}{We developed a coding schema}, \pc{huckvale2015af}{5}{Coding schema used to assess the privacy and security-related content of policy}. Each topic was coded \emph{addressed / absent}. Rights were coded afforded vs.\ denied (no right was ever mentioned only to be denied) \pc{huckvale2015}{4}{Assessment proceeded by systematically coding each}. Data items were coded per policy as \emph{will transmit / will not transmit / not addressed} \pc{huckvale2015}{4}{Individual data items were}. Coding was done by two researchers independently with custom annotation software, and all decisions were reviewed to consensus \pc{huckvale2015}{4}{Coding decisions, as well as any relevant policy text}. No agreement statistic is reported. The app co-developed by two authors was assessed by a non-involved author and checked by an external person \pc{huckvale2015}{12}{One of the apps included in this study, SiKL for iOS}.

\cf{Analysis} Descriptive; the unit is the platform-independent app. Fisher's exact test for policy availability by platform and payment model, with a sensitivity analysis restricted to data-collecting apps to address confounding by information-only apps \pc{huckvale2015}{5}{exact test to calculate the two-tailed probability of an association}.

\cf{Key findings} 89\,\% transmitted data, and no app encrypted local storage. 66\,\% of apps sending identifiers did so unencrypted. 67\,\% had some policy, but only 28\,\% had one in-app and only 10\,\% of web policies mentioned the app \pc{huckvale2015}{7}{apps had some form of privacy policy (Table 3)}, \pc{huckvale2015}{8}{of these web disclosures specifically referred to the app}. Policies addressed retention in only 17\,\% \pc{huckvale2015}{8}{fewer than one-fifth (17}. No app transmitted what its policy said it would not, but 78\,\% transmitted items the policy did not address \pc{huckvale2015}{8}{However, while no apps transmitted information where a}. Permissions were requested without a visible purpose (e.g.\ geolocation in 11 of 16 cases) \pc{huckvale2015}{9}{a clear reason, such as a mapping}.

\cf{Limitations stated} A snapshot of one library at one time. No access to server-side or organisational handling, so the policy is an imperfect proxy for actual behaviour \pc{huckvale2015}{11}{Our approach, instead, relied on the degree}.

\cf{Use for this thesis} Strong models for: a \emph{frozen software state} during evaluation; a two-phase walkthrough (exhaustive UI sweep, then realistic use); simulated plausible data with test flags; a search protocol for privacy text beyond the policy (terms, disclaimers, store marketing); a three-valued per-item coding (will / will not / not addressed); and an \emph{afforded vs.\ denied} distinction for rights statements. The latter is directly useful for coding the \emph{stance} of protection statements. Traffic and server testing are out of scope.
\end{papercard}


\begin{papercard}{huckvale2015af}
\cf{What it is} The 20-page Additional file 1 of Huckvale et al.\ (2015): data-type taxonomy (Table AF1), an explanatory figure on intercepting encrypted traffic (AF2), screenshots of the custom review and annotation software (AF3, AF5), the policy coding schema (AF4), and tables of excluded (AF6) and included (AF7) apps and of data transmissions (AF8). It documents the instrument behind the main paper, which is what makes the study reproducible in principle.

\cf{Data-type taxonomy (AF1)} Seven groups: \emph{strong personal identifiers} (``personal information that alone may uniquely identify an individual or allow identity to be compromised'') \pc{huckvale2015af}{1}{Personal information that alone may uniquely identify an individual or allow identity to be compromised}, \emph{weaker identifiers} that may be combined or linked \pc{huckvale2015af}{1}{Personal information that may be combined or linked to other data to uniquely identify an individual}, \emph{health-related information} (disability, genetic, HIV, IVF, measured parameters, medical history, medications, mental health status, substance use) \pc{huckvale2015af}{1}{Sensitive personal information relating to health status and medical history}, \emph{other sensitive information} (employment, ethnicity, political, religious, sexuality), \emph{other user-generated data}, \emph{analytics data} and \emph{protocol-standard data} (device characteristics, IP address), defined as data ``sent routinely as part of network communications over which app has limited control'' \pc{huckvale2015af}{2}{Data sent routinely as part of network communications over which app has limited control}. Geolocation and gender are listed as weaker identifiers, not as health data.

\cf{Interception and review tools} Figure AF2 explains how a man-in-the-middle set-up can read encrypted traffic by impersonating the device to the server and presenting a bogus certificate, with the caveat that device settings must be altered so that it accepts the certificate and that ``technical details have been simplified'' \pc{huckvale2015af}{3}{Technical details have been simplified to communicate key principles}. The custom software paired each outgoing request with its response so that a reviewer could annotate the purpose of the message and the data types it carried \pc{huckvale2015af}{4}{The overall purpose of the message as well as transmission of specific data types could then be annotated using the tool}; a second tool supported step-by-step annotation of policy text against the schema \pc{huckvale2015af}{7}{Software supported a step-by-step process of review for each policy document}.

\cf{Policy coding schema (AF4)} Four domains: \emph{uses of data} (primary and secondary uses, sending data to developer services, advertisers or marketers, analytics or research, while loading content, anonymous uses only) \pc{huckvale2015af}{5}{Primary uses of collected data (e.g. administering accounts}; \emph{technical concerns} (security arrangements, retention, inherent risks of mobile and internet security, cookies); \emph{user rights} (opt-out, consequences of not sharing, subject access requests, editing and deletion, complaints, vulnerable users and minors) \pc{huckvale2015af}{5}{Special procedures for handling data for vulnerable users and minors}; \emph{administrative details} (data controller, governing jurisdiction, processing jurisdictions, date of policy, next review, procedure for changing terms) \pc{huckvale2015af}{5}{Jurisdictions under which transmitted data will be processed}.

\cf{Exclusion and inclusion tables (AF6, AF7)} Exclusion reasons are listed app by app (unable to log in, not available for download, would not start, cost above the threshold, duplicate), for example an app that on iOS could not log in and on Android would not start \pc{huckvale2015af}{8}{Unable to log in (iOS). Would not start (Android)}. A footnote clarifies that the iOS versions of two apps were included because only their Android versions failed, so they are not counted among the excluded apps \pc{huckvale2015af}{9}{iOS versions of these apps were not affected by installation or start-up problems and were included in the study}. The included-app table gives name, platform and version, developer, cost and feature categories (information, health promotion, diary or personal health record, self-management, self-assessment, symptom checker, social network, assistive technology, service directory, reminders and others) \pc{huckvale2015af}{10}{Characteristics of included apps}.

\cf{Transmission summary (AF8)} 70 apps transmitted data; 63 (90\,\%) sent data to third parties, 50 (79\,\%) of these to destinations outside the UK \pc{huckvale2015af}{20}{Third parties 63 (90\%) 50 (79\%)}; analytics collection 43 (61\,\%) with 42 (98\,\%) outside the UK \pc{huckvale2015af}{20}{Analytics data collection 43 (61\%) 42 (98\%)}; marketing or advertising 14 (20\,\%) \pc{huckvale2015af}{20}{Marketing or advertising 14 (20\%) 9 (64\%)}. \emph{Check before quoting}: three cells in the text layer do not compute (``78 (305\%)'' for developer-controlled services; ``1 (10000\%)'' for research data collection to third parties; ``2 (10\%)'' for third-party account services when 2 of 70 is about 3\,\%), so use the main article for exact figures.

\cf{Use for this thesis} (1) AF1 is a ready taxonomy for the ``data type'' code in the policy codebook and for aligning with store labels; note that location and gender sit with weaker identifiers while mental health and substance use sit in the health group. (2) AF4 provides four domains and their topics as a cross-check of the thesis codebook's coverage; the ``date of policy'' and ``procedure for changing terms'' topics are easy to forget. (3) AF6 is a good pattern for a transparent exclusion table: every excluded app with a reason and a note when a platform version was kept. (4) The thesis does not intercept traffic, so AF2 and AF3 are context only; say so.
\end{papercard}
